\section{Related Literature}
\label{sec:literature}

\subsection{Connectivity, complementary inputs and adjustment}

The economic premise is that connectivity complements some productive tasks while others can be performed locally. \citet{hjort2019} study the arrival of submarine Internet cables and terrestrial network access in twelve African countries using difference-in-differences. They find higher employment, with firm evidence linking the expansion to entry, productivity and exports. Their setting concerns the adoption of a better technology over time; it does not imply that suddenly removing access has an equal and opposite effect. Existing users may have accumulated complementary skills and organizations, and a short outage leaves little time to reorganize production. Bangladesh's episode therefore provides evidence on the short-run dependence of an established digital activity on connectivity, rather than the long-run return to expanding broadband.

\citet{akerman2015} use a public broadband rollout in Norway as a source of variation in firms' adoption. Their evidence supports complementarity with skilled workers performing nonroutine abstract tasks and substitution for unskilled routine work. The relevant implication here is task heterogeneity: local editing may remain feasible without Internet access, while publishing, reviewing and coordinating changes depend directly on online services. It would be inappropriate to infer workers' skill or occupation from a repository label alone. Instead, I distinguish observable tasks and study whether users associated with different repository domains exhibit different public-activity gaps.

A related literature asks how firms adapt when an essential infrastructure input fails. \citet{allcott2016} instrument Indian manufacturers' electricity shortages with hydroelectric supply variation. Their estimates distinguish revenue and producer-surplus losses from smaller productivity losses, in part because inputs can be stored; access to generators also changes exposure. The parallel for software is the possibility of storing local work and uploading it later. The analogy is a mechanism hypothesis, not an assumption that software firms and factories adjust identically. I examine it by separating commit-author timestamps from publication timestamps and reporting delayed output both as a share of observed commits and relative to normal output levels.

\subsection{Restrictions, networks and public software activity}

Connectivity restrictions may affect people who retain access through their collaborators. \citet{zhang2011} exploit the blocking of Chinese Wikipedia and document reduced contributions among users who were not themselves blocked. This evidence on participation incentives makes cross-border spillovers relevant to the present design: a donor-country developer working with Bangladeshi collaborators need not be wholly untreated. A national synthetic control is credible only to the extent that such spillovers and donor-specific disturbances do not explain the counterfactual path. \citet{aceto2018} survey the measurement and causes of Internet outages, while \citet{schachenhofer2023} analyze organizational information flows. These studies motivate the distinction between connectivity, coordination and final output, but do not by themselves identify the effect of Bangladesh's July episode. Cross-country shutdown-cost models such as \citet*{tagat2024} and the regional application in \citet{gunsilius2023} address broader economic activity; this paper's contribution is a directly observed, high-frequency public software margin with a transparent denominator.

GitHub traces are useful because they record actions and timing, but their economic interpretation depends on the measure. \citet{mcdermott2021} use activity from millions of accounts to study labor reallocation during COVID-19, including changes toward weekends and nonstandard working times. Their approach illustrates why an aggregate count can change because of timing rather than a proportional change in labor supplied. \citet{wachs2022} map the geography of open-source software and \citet{birkholz2026} study international software-production networks. These applications motivate country comparisons while also making self-reported locations and cross-border collaboration central measurement concerns. \citet{nagle2019} relates firms' use of open-source software to productivity; the economic relevance of open source does not, however, provide a conversion factor from an individual GitHub event to value added. Likewise, \citet{hu2023} study pandemic lockdowns in China, a disruption whose effects on time allocation and online work need not resemble loss of connectivity.

The closest measurement precedent is \citet{kreitmeir2024}, who combine Italy's temporary ChatGPT ban with daily GitHub output for Italian and other European users in a difference-in-differences design. They report heterogeneous responses by experience and task type. Their conversion of sparse daily actions into indicators motivates my extensive-margin outcome, but the treatments differ: restricting one tool leaves Internet access available, whereas a network outage impedes publication and collaboration themselves. Accordingly, I separately define public activity, shipping and coordination, rather than treating them as interchangeable measures of total productivity. The commit-timing analysis further asks whether apparently missing publication is observed later. It cannot recover work that never becomes public or validate the physical time and location of authorship.

\subsection{Comparative designs and the contribution of this study}

\citet{abadie2003} introduce a synthetic comparison in their study of conflict in the Basque Country. \citet*{abadie2010} formalize the method, and \citet{abadie2015,abadie2021} emphasize the relationship between donor choice, predictor balance and the credibility of extrapolating the pre-treatment relationship. This framework is appropriate to a country-level episode, but it does not remove concurrent unrest or compensate automatically for poor fit. \citet{ferman2021} show why imperfect pre-treatment fit deserves explicit attention. I therefore report the optimization problem, all candidate donors, lag specifications, holdout fit, placebo comparisons and donor exclusions. The fitted weights are empirical solutions, not claims that the positive-weight countries are structurally identical to Bangladesh.

Alternative estimators answer related questions under additional choices. Synthetic difference-in-differences adds time reweighting \citep{arkhangelsky2021}; here that reweighting emphasizes the unusually active final pre-shutdown weeks, so its larger estimate is presented as sensitivity evidence. Regression discontinuity in time requires care about serial correlation, functional form and coincident shocks \citep{hausman2018}; its details are relegated to the appendix. The central contribution is the combination of a transparent national counterfactual, task-specific public outcomes, fixed pre-treatment repository-domain heterogeneity, and a separately defined residual-author cohort. These components distinguish disruption from observed adaptation while preserving the limits of what public traces can establish.
